% =====================================================================
%  ZUSAMMENFASSUNG UND DISKUSSION  (Abschnitt 1.5 der Anleitung)
%  Ergebnisse inkl. Fehlerangaben und Literaturwerten in derselben
%  Einheit; anschließend Diskussion: statistische vs. systematische
%  Fehlerquellen, Vergleich mit Literatur.
% =====================================================================
\section{Zusammenfassung und Diskussion}
\label{sec:diskussion}

\subsection{Zusammenfassung der Ergebnisse}
Alle berechneten Ergebnisse inkl. Fehlerangaben und Literaturwerten,
in derselben Einheit (ggf. umrechnen):

\begin{table}[H]
  \centering
  \caption{Ergebnisse im Vergleich mit Literaturwerten.}
  \label{tab:ergebnisse}
  \begin{tabular}{llll}
    \toprule
    Größe & eigener Wert & Literaturwert & Quelle \\
    \midrule
    $v$ & $(2{,}1 \pm 0{,}1)\,\si{\milli\meter\per\second}$
        & $2{,}0\,\si{\milli\meter\per\second}$ & \cite{Example2023} \\
    \bottomrule
  \end{tabular}
\end{table}

\subsection{Diskussion}
Bewertung der Ergebnisse im Vergleich mit den Literaturwerten:

\begin{itemize}
  \item \textbf{Systematische Fehler} (Abschnitt 6.1 der Anleitung):
        reproduzierbare Abweichung in eine Richtung, z. B. falsch
        geeichte Messinstrumente, Kalibrierfehler. Wiederholung
        verbessert die Präzision nicht.
  \item \textbf{Statistische Fehler} (Abschnitt 6.2): Schwankungen
        äußerer Bedingungen (Temperatur, Luftdruck), Ablesegenauigkeit,
        Toleranzen der Experimentatoren. Wiederholung verbessert die
        Genauigkeit (siehe \cref{eq:fehler_mittelwert}).
  \item Grundsatz: ehrliche Abschätzung; im Zweifel Fehler lieber
        überschätzen als unterschätzen.
\end{itemize}
